% !TeX root = ../概率论与数理统计笔记.tex
% 由主文档通过 \input 载入；请编译主文档。
\setcounter{chapter}{19} % 保留原书章号。
\chapter{随机模拟与统计计算（选读）}
\label{chap:ps-computation}

本章用算法近似概率量和统计推断结果，建立理论对象、计算过程与误差评价之间的联系。
先修内容为极限定理与统计推断；马尔可夫链方法另需第\ref{chap:ps-processes}章。
每种算法的正文均明确输入、输出、初始化、更新、停止规则及适用条件。

\section{随机数与分布抽样}
\label{sec:ps-random-generation}
区分数学上的独立随机样本与计算机\emph{伪随机数}（pseudorandom numbers），记录种子和生成方式以便复现实验。
讲解逆分布函数抽样与拒绝抽样；后者要求可计算的目标与提议密度关系，以及有效的包络常数。
说明实际算法输出的是有限精度数值，理论分布性质需要结合实现方式理解。

\section{蒙特卡罗积分与误差估计}
\label{sec:ps-monte-carlo}
介绍\emph{蒙特卡罗方法}（Monte Carlo method），将积分写成期望，再以独立样本均值估计。
有限绝对一阶矩用于大数定律保证；有限非零方差用于通常的中心极限定理误差近似。
以预设模拟次数或有理论依据的序贯规则停止，不能不加说明地反复检查普通固定样本区间后自行停止。

\section{重要性抽样与方差缩减}
\label{sec:ps-variance-reduction}
介绍\emph{重要性抽样}（importance sampling）、控制变量与对偶变量的基本思想。
重要性抽样要求提议分布覆盖目标积分的有关支撑，并检查加权变量的可积性与二阶矩。
区分普通加权估计与自归一化估计的有限样本性质；比较方法时同时报告精度与计算成本。

\section{自助法与置换计算}
\label{sec:ps-bootstrap}
介绍\emph{自助法}（bootstrap），从经验分布重复抽样近似统计量的抽样分布，计算标准误与区间。
普通非参数自助法以适当的独立同分布抽样和统计量正则性为基础；极值、边界或依赖数据需要另行论证。
区分自助重抽样与原假设下的置换，两者的目标和合法性依据不同；有限次重复还引入模拟误差。

\section{马尔可夫链蒙特卡罗}
\label{sec:ps-mcmc}
介绍\emph{马尔可夫链蒙特卡罗}（Markov chain Monte Carlo, MCMC），先在有限状态情形构造以目标分布为平稳分布的链。
给出梅特罗波利斯--黑斯廷斯算法（Metropolis--Hastings algorithm）与吉布斯抽样（Gibbs sampling）的更新思想。
分别讨论不变分布、不可约性、非周期性与时间平均收敛所需条件；连续状态推广需补充相应理论。
相关输出不能直接当作独立样本，诊断图与有效样本量只提供计算证据，不能单独证明已经收敛。

\section{例题、习题与学习检查}
规划用模拟核对概率计算、比较两种积分估计方差、模拟检验功效与区间覆盖率，以及用简单离散目标分布检查马尔可夫链抽样。
每个实验同时报告模型、重复次数、统计量、随机性来源和模拟标准误；与可解析的小模型作交叉核对。
学习检查区分模型误差、抽样误差、蒙特卡罗误差与数值优化误差。




